\begin{frame}{멀티헤드 공식}
  모델 차원 $d_{model}$ 을 $h$개의 헤드에 나눠, $i$번째 헤드는
  $d_k = d_{model}/h$ 차원의 Q, K, V 위에서 12.2절의 어텐션을
  계산:
  \[
  \text{MultiHead}(Q,K,V) =
  \text{Concat}(\text{head}_1, \dots, \text{head}_h)\, W_O,
  \qquad
  \text{head}_i = \text{Attention}\!\left(QW_i^Q,\; KW_i^K,\;
  VW_i^V\right)
  \]
  \begin{itemize}
    \item $W_i^Q, W_i^K \in \mathbb{R}^{d_{model} \times d_k}$,
      $W_i^V \in \mathbb{R}^{d_{model} \times d_v}$:
      \textbf{헤드별 투영 행렬} — 각 헤드가 ``무엇과 무엇을
      비교할지''(Q·K)와 ``무엇을 전달할지''(V)를 \emph{각자}
      학습(12.1절의 Q/K/V 삼분 구조가 \textbf{헤드마다 하나씩}).
    \item $W_O \in \mathbb{R}^{h d_v \times d_{model}}$:
      이어붙인 결과를 다시 $d_{model}$ 차원으로 — ``각 헤드의
      정보를 어떻게 섞어 다음 층에 줄지''를 결정하는
      \textbf{최종 선형변환}.
  \end{itemize}
\end{frame}
